\documentclass[11pt,a4paper]{article}
\usepackage[T1]{fontenc}
\usepackage{lmodern,amsmath,amssymb,microtype,booktabs}
\usepackage[margin=25mm]{geometry}
\usepackage[colorlinks=true,urlcolor=blue,linkcolor=blue]{hyperref}
\newcommand{\HH}{\mathcal H}\newcommand{\MM}{\mathcal M}
\title{Regulated Fock fields and vacuum returns\\in the dS$_3$ observer model}
\author{Alan Day \and Chengheng\thanks{AI research assistant (OpenAI).}}
\date{Working research note --- 31 October 2026}
\begin{document}\maketitle
\begin{abstract}
We construct all finite-particle creation blocks from normalized
$\mathrm{SL}(2,\mathbb C)$ Clebsch--Gordan synthesis and identify their
weak group-averaged action on the compatible full multiplicity spaces.
Uniform one-step bounds imply square-root particle-number bounds and
essential self-adjointness on the finite-particle domain.
We distinguish the adjoint completion of half-period creation from
real Gaussian smearing of the original local field. Restoring the
vacuum-return path completes the four-field vacuum pairing. The
Gaussian heavy-observer limit has thermal detailed balance and a
generally nonzero free-field commutator, even though the two-creation
ordering difference vanishes. These are regulated $G_N=0$ results,
not a finite-gravity holographic dual.
\end{abstract}
\section{Two regulators with different meanings}
Take $\ell=1$, $G_N=0$, proper $SO^+(3,1)$, two distinguishable BD
scalar observers with labels at least $s>0$, and one free KG matter
species $p>0$. The local coupling is one; clock energy is
$h(a)=\sqrt{1+a^2}$. Keep the previously derived emission coefficient
$E_{abp}$ and its full angular CG tensor \cite{prior,kl}.
The finite-mode kernel is the spacetime integral of observer modes
and the matter creation mode.

Consider two emission multipliers:
\begin{align}
 m_h(a,b;t)&=e^{\pi[h(a)-h(b)]/2}e^{-it[h(a)-h(b)]},\nonumber\\
 m_\sigma(a,b;t)&=
 e^{-\sigma^2[h(a)-h(b)]^2/2}e^{-it[h(a)-h(b)]},\quad \sigma>0.
 \label{reg}
\end{align}
For each choice complete creation by its Hilbert adjoint:
$\Phi_j(t)=C_j(t)+C_j(t)^*$, $j=h,\sigma$.
The Gaussian choice is real time smearing of both parts of the original
Hermitian local field, initially understood on finite-mode weak kernels.
The $h$ choice completes the preceding half-period creation block.
Its annihilation factor is the adjoint factor; it is not obtained by
shifting the entire unregulated field to the same complex time.
We do not identify these observables.

Full representations and all internal spins are retained. A scalar
principal series contains all its $SO(3)$ angular modes.
No angular zero-mode truncation is made.

\newpage
\section{The bounded one-step local intertwiner}
Let $P_a$ be the scalar principal series and
$\HH_O=\int_{a\ge s}^{\oplus}da\,P_a$. The preceding normalization gives
\[
 \rho(b)=\frac{2b^2}{\pi},\quad
 K_3(a,p,b)=\frac{\pi}{8D_{apb}}
 \prod_{\nu=a,p,b}\frac{\sinh\pi\nu}{\nu},\quad
 w_a(b)=\rho(b)K_3(a,p,b).
\]
Here $D_{apb}$ is the vertex note's product of four cosh factors.
Normalized scalar-channel CG synthesis \cite{ds} is an isometry
\[
 J_a:\int_{b>0}^{\oplus}db\,P_b\longrightarrow P_a\otimes P_p .
\]
Normalize the unscaled CG tensor to have $000$ coefficient one.
The synthesis kernel is then $\sqrt{w_a(b)}C^\dagger_{a,p;b}$.
The actual emission kernel determines
\begin{equation}
 (T_j(t)\xi)(a)=
 J_a\!\left[\frac{E_{abp}m_j(a,b;t)}{\sqrt{w_a(b)}}\xi(b)\right],
 \qquad \xi(b)=0\quad(b<s).
 \label{T}
\end{equation}
This is $L^2$ spectral synthesis, not a statement that a fixed-$b$
CG column is square summable. Its finite-mode weak matrix elements
are the original local ones, with multiplier \eqref{reg}.

For the balanced choice the squared coefficient is
\[
 R_h(a|b)=\frac{|E_{abp}|^2e^{\pi[h(a)-h(b)]}}{\rho(b)K_3(a,p,b)}
         =\frac{g(a)}{g(b)}\widehat R_p(a,b),
\]
\[
 g(x)=\frac{x e^{\pi[h(x)-x]}}{1+e^{-2\pi x}},\quad
 0\le\widehat R_p(a,b)\le k_p(a-b),\quad
 k_p(u)=\frac1{8\pi(\cosh\pi u+\cosh\pi p)} .
\]
The bound $g(a)/g(b)\le B_s(1+|a-b|/s)$,
$B_s=2e^{\pi/(\sqrt{1+s^2}+s)}$, gives
\begin{equation}
 \sup_{b\ge s}\int_s^\infty R_h(a|b)\,da
 \le {\cal B}_h=B_s\left(\frac1{2\pi^2}+\frac1{2s\pi^3}\right).
 \label{B}
\end{equation}
CG isometry proves $\|T_h(t)\|^2\le\mathcal B_h$.
For Gaussian smearing, completing the square gives
\begin{equation}
 R_\sigma=e^{-\sigma^2\Delta h^2-\pi\Delta h}R_h
 \le\Lambda_\sigma R_h,\qquad
 \Lambda_\sigma=e^{\pi^2/(4\sigma^2)},\quad
 \mathcal B_\sigma=\Lambda_\sigma\mathcal B_h ,
 \label{GaussianB}
\end{equation}
where $\Delta h=h(a)-h(b)$. Thus $\|T_\sigma(t)\|^2\le\mathcal B_\sigma$.
These are bounded intertwiners: coefficients depend only on
representation labels. Normalization comes from the local vertex,
not covariance alone. Bounds are not uniform as $s\downarrow0$
or $\sigma\downarrow0$.

\newpage
\section{All particle numbers and their physical completion}
Put $W_N=\HH_O\otimes\operatorname{Sym}^N P_p$; the other observer is
the spectator. The full principal-series tensor-product theorem
\cite{ds,my} includes allowed integer spins and continuous parameters.
In this dimension no discrete tempered series is added.
Iterating the full unitary CG transform puts all previous internal
paths into multiplicity spaces. Bose symmetrization commutes with
the group and projects these multiplicities; it can mix paths and
does not discard nonzero-spin internal channels.

Only a scalar total representation pairs with the scalar spectator.
It contains one $K=SO(3)$ invariant vector, while nonzero-spin
principal series contain none. Thus
\begin{equation}
 W_N^K\simeq\int_{c>0}^{\oplus}dc\,\MM_{N,c}.
 \label{WN}
\end{equation}
All internal spins remain in $\MM_{N,c}$.
For $N=0$, $\MM_{0,c}=\mathbb C$ for $c\ge s$ and zero otherwise;
a fixed single observer has no extra independent total spectral label.
Use seeds
\[
 {\cal E}_N=\operatorname{span}\{|f,0\rangle_L\otimes v:
       f\in C_c^\infty((s,\infty)),\ v\in W_N^K\}.
\]
The spherical function, radial Haar measure and uniform bound are
\[
 \phi_c(r)=\frac{\sin cr}{c\sinh r},\qquad
 \sinh^2r\,dr,\qquad |\phi_c(r)|\le\frac r{\sinh r}.
\]
For $F=\overline{f_1}f_2$, integrate the left mass first.
The factor $\int F(c)\sin(cr)/c\,dc$ decreases faster than any
inverse power of $r$ by integration by parts.
The right coefficient is bounded by $(r/\sinh r)\|v\|\|w\|$.
After the Haar density the product is bounded at infinity by
$C_n r(1+r)^{-n}$, integrable for $n>2$.
It is regular at zero. Absolute group-average convergence follows
for every $\mathcal E_N$, including mass-smeared $N=0$.

Sine-transform inversion, initially for compact smooth spectral data,
then extended with the uniform majorant, yields
\begin{align}
 Q_N(f\otimes v,k\otimes w)
 &=\int_{c\ge s}\frac{dc}{\rho(c)}
      \overline{f(c)}k(c)\langle v(c),w(c)\rangle,\nonumber\\
 A_N(f\otimes v)(c)&=\rho(c)^{-1}f(c)v(c),&
 \HH_N&=\int_{c\ge s}^{\oplus}\rho(c)\,dc\,\MM_{N,c}.
 \label{GA}
\end{align}
This does not take the absolute Haar integral of two fixed-label
spherical functions. Fubini is used after left mass integration,
with the right spectral $L^1$ pairing and uniform spherical bound.
The image is dense: for compact $c$-supported $F(c)\in\HH_N$,
choose $f=1$ on its support and $v=\rho F$ there. Then $v\in W_N^K$.
In addition
\[
 \|A_N(f\otimes v)\|\le
 \left(\operatorname*{ess\,sup}|f|^2/\rho\right)^{1/2}\|v\|.
\]
Finite-angular-mode, compact actual observer-mass approximations
to $v$, projected onto total $K$ invariance, give the same completion.

\newpage
\section{Weak descent and self-adjoint fields}
Ordinary Fock symmetrization defines the kinematic maps
\[
 {\cal C}_{j,N}(t)=\sqrt{N+1}
 (1_O\otimes P_{\mathrm{Sym}^{N+1}})
 (T_j(t)\otimes1),\qquad
 \|{\cal C}_{j,N}(t)\|\le\sqrt{(N+1)\mathcal B_j}.
\]
They are intertwiners. In the full irreducible decomposition their
fiber maps $c_{j,N}(c;t)$ have the same essential uniform bound.
The maps and their adjoints preserve total $K$-invariant subspaces.

Take a common compact group cutoff $0\le\chi_R\le1$, inverse symmetric,
tending to one, and $\eta_R=\int\chi_R(g)U(g)\,dg$.
At finite $R$, bounded intertwining commutes with this Bochner integral.
The local weak creation matrix element is consequently
\begin{align}
 \lim_R\langle\Phi,(1_L\otimes{\cal C}_{j,N})\eta_R\Psi\rangle_{\rm kin}
 &=Q_{N+1}(\Phi,(1_L\otimes{\cal C}_{j,N})\Psi)\nonumber\\
 &=\int_{c\ge s}\rho(c)\,dc\,
 \langle A_{N+1}\Phi,c_{j,N}(c;t)A_N\Psi\rangle .
 \label{weak}
\end{align}
The enlarged seed spaces are preserved, so the absolute-convergence
estimate applies and makes the limit cutoff independent.
The expression descends through both null spaces; density gives a
unique bounded physical block $C_{j,N}(t)$.
Using the same inverse-symmetric cutoff identifies the reverse weak
block with its adjoint. This is the test-domain step needed for descent
\cite{raq}, rather than a presumption based on covariance.

The initial finite-mode kernels can also be approximated by finite-rank
matrices before group averaging. At each fixed group cutoff, strong
convergence and the uniform block bound recover the intertwiner.
Only then is the group cutoff removed. No exchange of these two
limits, or action on arbitrary unsmeared invariant distributions,
is needed for \eqref{weak}.

Let $\HH=\bigoplus_{N\ge0}\HH_N$ and
$\mathcal D_{\rm fin}=\bigoplus_{N\ge0}^{\rm alg}\HH_N$.
Each $\Phi_j(t)=C_j(t)+C_j(t)^*$ is symmetric on this invariant dense
domain. For $\psi$ supported through particle number $r$,
\[
 \|\Phi_j(t)^k\psi\|\le
 (2\sqrt{\mathcal B_j})^k\sqrt{\frac{(r+k)!}{r!}}\,\|\psi\|.
\]
The analytic-vector series converges for every argument by the ratio
test. Nelson's theorem \cite{nelson} proves essential self-adjointness.
Finite products at different real times also exist on this common
domain. The closed fields are unbounded, though each sector block
has the stated square-root bound.

This realizes the full compatible scalar multiplicity space, including
all internal channels. It does not establish maximality among
unspecified invariant-distribution domains, or a pointwise
position-space formula for every completed vector.
No unregulated point-field conclusion follows.

\newpage
\section{Vacuum returns and the exact four-field identity}
Let $\Omega_f$ have zero matter particles and
$\int_s^\infty\rho(c)|f(c)|^2dc=1$.
Define
\begin{equation}
 q_{j,c}(\tau)=\int_s^\infty da\,R_j(a|c)
             e^{-i\tau[h(a)-h(c)]},\qquad j=h,\sigma .
 \label{return}
\end{equation}
On the one-dimensional vacuum multiplicity fiber,
$C_{j,0}(t_1)^*C_{j,0}(t_2)=q_{j,c}(t_2-t_1)$.
This is diagonal in the conserved spectator label $c$, not in the
actual right observer label in higher particle sectors.
The exact vacuum four-field expectation is
\begin{align}
 &\langle\Omega_f,\Phi_j(t_1)\Phi_j(t_2)
                        \Phi_j(t_3)\Phi_j(t_4)\Omega_f\rangle\nonumber\\
 &\quad=\int_s^\infty\rho(c)|f(c)|^2
       q_{j,c}(t_2-t_1)q_{j,c}(t_4-t_3)\,dc\nonumber\\
 &\qquad+
 \langle C_{j,1}(t_2)C_{j,0}(t_1)\Omega_f,\,
         C_{j,1}(t_3)C_{j,0}(t_4)\Omega_f\rangle .
 \label{four}
\end{align}
Read from right to left, exactly two particle-number paths survive:
$0\to1\to0\to1\to0$ and $0\to1\to2\to1\to0$.
The second term contains direct and exchange contractions from Bose
symmetrization. The first integral is not a product of separately
mass-averaged two-point functions: both returns share the same $c$.

The squared commutator norm decomposes exactly as
\begin{equation}
 \|[\Phi_j(t),\Phi_j(0)]\Omega_f\|^2
 =D_{B,j,f}(t)+4\int_s^\infty\rho(c)|f(c)|^2
                           [\Im q_{j,c}(t)]^2\,dc .
 \label{comm}
\end{equation}
The terms occupy orthogonal two-particle and vacuum sectors.
The vacuum contribution was absent from the creation-only observable.
At $t=0$ both vanish exactly.

Truncating the balanced transfer $u=a-c$ at $|u|\le D$ gives
\begin{equation}
 \epsilon_1(D)=\frac{B_se^{-\pi D}}{2\pi}
 \left[\frac1\pi+\frac1s\left(\frac D\pi+\frac1{\pi^2}\right)\right],
 \label{onetail}
\end{equation}
uniformly for $c\ge s$ and real $\tau$.
For Gaussian smearing multiply by $\Lambda_\sigma$.
This integrates the one-step exponential majorant on both tails,
including its harmless extension beyond the physical threshold.
The return functions are bounded by $\mathcal B_j$, controlling the
mass averages in \eqref{four} and \eqref{comm}.
Analytic tails exclude quadrature and arithmetic errors.

\newpage
\section{Heavy-observer spectra and thermal detailed balance}
Fix $s,p,\sigma>0$ and condition $c=M\to\infty$.
The limiting densities and transforms are
\begin{align}
 k_h(u)&=k_p(u),&
 L_h(\tau)&=\frac{\sin(p\tau)}
                    {4\pi\sinh(\pi p)\sinh\tau},\nonumber\\
 k_\sigma(u)&=e^{-\sigma^2u^2-\pi u}k_p(u),&
 L_\sigma(\tau)&=\int_{\mathbb R}k_\sigma(u)e^{-i\tau u}\,du .
 \label{limits}
\end{align}
Here $L_h(0)=p/[4\pi\sinh(\pi p)]$.
Dominated convergence follows from the one-step majorant and
\eqref{GaussianB}, uniformly on compact real times. In either case
\[
 q_{j,M}(\tau)=L_j(\tau)+\frac{iL_j'(\tau)}M+O_T(M^{-2}).
\]
On $|u|\le M/2$ use $g(M+u)/g(M)=1+u/M+
O((|u|+u^2)/M^2)$ and
$h(M+u)-h(M)=u+O(|u|/M^2)$.
The balanced-kernel error is exponential; the omitted region has an
exponentially decreasing polynomial majorant.
The Gaussian multiplier and its derivatives are bounded after
completing the square, so the same argument applies.

The limiting observables differ: $k_h$ is even and $L_h$ is real and
even, with no finite-temperature detailed balance at nonzero frequency.
The Gaussian local field instead has
\begin{equation}
 k_\sigma(-u)=e^{2\pi u}k_\sigma(u),\qquad
 L_\sigma(\tau+2\pi i)=L_\sigma(-\tau).
 \label{KMS}
\end{equation}
The Gaussian Fourier integral is entire; changing the integration
variable proves the identity. Our convention is $\tau=t_2-t_1$.
Thus $W(t)=\langle\Phi(t)\Phi(0)\rangle=L_\sigma(-t)$ has the usual
lower-strip KMS convention. Exact finite-mass thermality is not claimed.

The previous compact crossing theorem applies to arbitrary four times
in a fixed compact set. For $h$, the $C_r^2$ physical weight simply
acquires a general linear combination of energy phases.
For $\sigma$, each emission differs from the balanced one by a smooth
multiplier of modulus at most $\sqrt{\Lambda_\sigma}$.
A two-emission chain is bounded by $\Lambda_\sigma$ times its
balanced counterpart. Its norm and deleted squared norm bounds,
and the crossed Hilbert tail, gain $\Lambda_\sigma^2$.
Thus the fixed-compact limit followed by cutoff removal still applies.

At the ridge $r=0$, the direct phase is
$e^{-iv(t_3-t_2)-iu(t_4-t_1)}$ and the exchange phase is
$e^{-iu(t_3-t_1)-iv(t_4-t_2)}$.
The ridge spectral weight is $k_j(u)k_j(v)$.
These fix the two contractions and their time order, without assuming
finite-mass Wick factorization of the observer matrices.

\newpage
\section{Full four-point limit and its commutator}
Both two-creation contractions and the vacuum return combine to give
\begin{align}
 G_{4,j}^{\infty}(t_1,t_2,t_3,t_4)
 &=L_j(t_2-t_1)L_j(t_4-t_3)\nonumber\\
 &\quad+L_j(t_3-t_2)L_j(t_4-t_1)\nonumber\\
 &\quad+L_j(t_3-t_1)L_j(t_4-t_2).
 \label{Wick}
\end{align}
This is the standard three-contraction Gaussian four-point baseline,
derived here from the specified observer construction.
Convergence is uniform for all four times in a fixed compact set.
It also holds for normalized probabilities on $[s,\infty)$ whose mass
escapes every fixed compact interval: split the average at $c=C$ and
use the uniform block bounds. A fixed nonmoving packet does not meet
this preparation hypothesis.

The two-creation difference obeys $D_{B,j,M}(t)\to0$ in both cases,
whereas the complete field gives
\begin{equation}
 \|[\Phi_j(t),\Phi_j(0)]\Omega_M\|^2
       \longrightarrow4[\Im L_j(t)]^2.
 \label{commfree}
\end{equation}
The balanced limit vanishes, consistently with its even spectrum.
The Gaussian local-field limit is generally nonzero: it is the
free-field commutator contribution.
Thus the earlier $D_B\to0$ cannot be promoted to unequal-time
commutativity of the thermal local field.

As an additional balanced-return diagnostic, second-order expansion gives
\begin{equation}
 q_{h,M}(t)=L_h(t)+\frac{iL_h'(t)}M
                -\frac{(t+i\pi)L_h'(t)}{2M^2}+O_T(M^{-3}).
 \label{second}
\end{equation}
The weight ratio contributes $-\pi u/(2M^2)$ and the energy phase
contributes $itu/(2M^2)$. The same inner/outer split controls the
integrated remainder. Its vacuum term starts as
$4L_h'(t)^2/M^2+O_T(M^{-3})$.
This is not a leading-rate claim for the total commutator, because
an optimized global mass-rate estimate for its two-creation part
is still unavailable.

The result is a regulated weak self-adjoint field on the stated full
multiplicity completion and its vacuum four-point limit.
We have not removed regulators, proved uniqueness among arbitrary
larger distributional spaces, or determined the observer algebra's
factor type. All-$N$ existence is not an all-correlation Gaussian-limit
theorem. Growing times, finite-$G_N$ dynamics and a microscopic
holographic identification require further work.

\newpage
\section{Numerical evidence and retained failures}
Balanced returns use 60 decimal digits, $p=0.6$, $s=0.2$,
$M=0.7,4,16,64,256$ and $t=0,0.5,1,2$.
All 20 cases retain $D=24$ and $D=32$ cutoffs, complex values and tails.
The previously omitted vacuum contribution at $t=1$ is:
\begin{center}
\begin{tabular}{rr}
\toprule
$M$ & $4[\Im q_{h,M}(1)]^2$\\
\midrule
0.7 & $2.4737552\,10^{-5}$\\
4 & $2.6287536\,10^{-6}$\\
16 & $3.4149846\,10^{-7}$\\
64 & $2.4951661\,10^{-8}$\\
256 & $1.6188726\,10^{-9}$\\
\bottomrule
\end{tabular}
\end{center}
These are empirical quadratures, not certified intervals.
Five direct $|E|^2/w$ assemblies match the simplified rate with maximum
relative difference about $2.02\,10^{-58}$.
An auxiliary model uses three observer levels and two oscillator modes,
with noncommuting complex observer matrices. Four block-norm checks,
four full four-field path decompositions and three orthogonal-sector
commutator checks pass. This tests finite Fock bookkeeping, not the
noncompact group-average theorem.

The driver saved all quadratures, then failed on JSON serialization
of a NumPy Boolean in the matrix results. A separate recovery reused
those quadratures unchanged and converted NumPy scalars explicitly.
The original script and raw outputs remain. The recovered diagnostic
set has 36 passing criteria.
An independent Decimal and vectorized 96/160-order Gauss audit does
not import the production return kernel. The largest reference
difference is about $4.34\,10^{-17}$; the largest order change is
about $8.25\,10^{-17}$. An initial structural SymPy equality rejected
equal algebraic expressions; the failure is retained. An exact
zero-polynomial check repairs it, giving 154 passing checks with no
data or tolerance changes. These checks do not certify total error.

Gaussian data use $p=0.6$, $s=0.2$, $\sigma=0.8$,
$M=4,16,64,256$ and the same four times.
Both cutoffs, limiting Fourier integrals and detailed-balance/KMS
comparisons are saved separately; all 24 criteria pass.
At $t=1$, the limiting vacuum contribution is approximately
$0.00411327832$, compared with $0.0040762936$ at $M=256$.
The largest KMS comparison residual is about $1.09\,10^{-62}$;
it is an empirical check of the analytic identity, not finite-mass KMS.
No new crossed $U$ or complete
four-field quadrature is claimed; \eqref{Wick} rests on the compact
and Hilbert-tail argument.
Dependencies, hashes and compilation records are under
\path{results/stage4/ds3-fock-vacuum/}.

\newpage
\section*{Sources and next research step}
The ingredients are established CG theory, group averaging,
bosonic symmetrization and the analytic-vector theorem.
This self-checked application connects them to the project's fixed
vertex, measure and regulators; no exhaustive priority claim is made.
Next compare these $G_N=0$ observables with finite-gravity insertions,
preserving the regulator and vacuum-return dictionary.
Free-limit correlations alone do not identify a factor type or
an entropy-counting dual.

\begin{thebibliography}{9}
\bibitem{prior} Alan Day and Chengheng, project notes
\emph{dS$_3$ local vertex}, \emph{Group pairing},
\emph{Two-matter sector}, \emph{Bose normalization},
\emph{Physical local weights} and \emph{Compact-transfer extension},
31 October 2026. The dS$_2$ all-$N$ and weak-descent notes are
comparisons, not substitutes for this dimension's proof.
\bibitem{ds} S. E. Derkachov and V. P. Spiridonov,
\href{https://arxiv.org/html/1711.07073v2}{arXiv:1711.07073v2},
Section 3.1, Eqs.\ (22), (29), (30), and Section 4.
Full spin sums and reflection representatives are essential.
\bibitem{my} D. Marolf and X. Yu,
\href{https://arxiv.org/html/2610.24073v1}
{On the Positivity of de Sitter Group Averaging},
Sections 5.1--5.2 and footnote 29.
Two observers without matter require mass packets, not a convergent
fixed-label two-particle Haar integral.
\bibitem{kl} D. K. Kolchmeyer and H. Liu,
\href{https://arxiv.org/html/2411.08090v1}
{Chaos and the Emergence of the Cosmological Horizon},
Sections 4--6, the observer and Wick-contraction framework in dS$_2$.
\bibitem{raq} D. Giulini and D. Marolf,
\href{https://arxiv.org/html/gr-qc/9902045v2}
{A Uniqueness Theorem for Constraint Quantization},
Section II. The weak-domain argument does not presume that an
unregulated local field preserves a rigging test domain.
\bibitem{nelson} E. Nelson, \emph{Analytic vectors},
Annals of Mathematics 70 (1959), 572--615.
The criterion was checked in the saved K.-H. Neeb,
\emph{On analytic vectors for unitary representations of infinite
dimensional Lie groups}, Ann.\ Inst.\ Fourier 61 (2011), Theorem 4.8.
\end{thebibliography}
\end{document}
